% form_tractatus.tex
% Form Tractatus - Sju kardinalproposisjonar for Form Follows Fitness
% STOLAR-prosjektet
% Spraak: Nynorsk

\documentclass[12pt, a4paper]{article}

\usepackage[utf8]{inputenc}
\usepackage[T1]{fontenc}
\usepackage[nynorsk]{babel}
\usepackage{mathpazo}
\usepackage{amsmath, amssymb}
\usepackage{xcolor}
\usepackage{hyperref}
\usepackage{setspace}
\usepackage{geometry}
\usepackage{csquotes}

\geometry{margin=3.5cm}
\onehalfspacing

\newcommand{\prop}[1]{\bigskip\noindent\textbf{#1}\quad}
\newcommand{\subprop}[1]{\medskip\noindent\textbf{#1}\quad}

\title{%
  {\huge\textbf{Traktat for form}}\\[1em]
  {\large Sju kardinalproposisjonar}\\[0.5em]
  {\normalsize Form Follows Fitness}
}
\author{Iver Raknes Finne}
\date{}

\begin{document}
\maketitle
\thispagestyle{empty}

\vspace{2em}

\begin{center}
\emph{Det deterministiske aksiomet har falle.\\Her byrjar tilpassinga.}
\end{center}

\vspace{3em}


% ================================================================
% PROPOSISJON 1
% ================================================================

\prop{1} Formverda er alt som er tilfelle i det $n$-dimensjonale formrommet.

\subprop{1.1} \textbf{Formrommet} er totaliteten av alle moglege geometriske, materielle og operasjonelle tilstandar eit system, ein organisme eller ein gjenstand kan innta. Ikkje alle posisjonar er busette; dei tome regionane fortel kva former som er fysisk, teknologisk eller kulturelt umoglege.

\subprop{1.12} Kvar realisert form er eit punkt i dette rommet. Kvar urealisert form er ein posisjon som anten er uoppdaga, umogeleg gjeve dei rådande materialaffordansane, eller ekskludert av dei gjeldande seleksjonstrykka.

\subprop{1.14} Det deterministiske aksiomet frå det tjuande hundreåret, postulatet om at form let seg ufeilbarleg utleie direkte frå funksjon, er eit logisk usant bilete av røynda.

\subprop{1.141} Aksiomet predikerer at objekt med same funksjon har same form. Prediksjonen er falsifisert: blant 2\,284 stolar med eksakt same invariante funksjon, å bere ein sitjande kropp, varierer gjennomsnittleg høgde frå 62 cm til 115 cm, materialpaletten frå 2 til 65 unike kategoriar, og estimert medianvekt frå 6,8 kg til 20,4 kg.

\subprop{1.142} Funksjonen er konstant; formvariasjonen er massiv. Aksiomet har ikkje eit einaste hundreår med empirisk støtte.

\subprop{1.2} Forma følgjer såleis ikkje funksjon; forma følgjer tilpassing (\emph{fitness}).


% ================================================================
% PROPOSISJON 2
% ================================================================

\prop{2} Det som er tilfelle, \emph{formfakta}, er eksistensen av lokale optimum i eit tilpassingslandskap.

\subprop{2.1} \textbf{Tilpassingslandskapet} er definert av tre element:
\begin{enumerate}
  \item Eit \emph{konfigurasjonsrom}: totaliteten av alle moglege formtilstandar.
  \item Ein \emph{naboskapsstruktur}: kva tilstandar som er tilgjengelege frå kvarandre, gjeve eksisterande produksjonsteknologiar og materialaffordansar.
  \item Ein \emph{tilpassingsfunksjon}: ein samansett verdi der funksjonell yting, estetisk aksept, produksjonskostnad, materialtilgang og kulturell meining alle bidreg.
\end{enumerate}

\subprop{2.12} Landskapet er djupt ruglete. Det har fleire lokale optimum, og formgjeving er navigasjon i dette terrenget. Kompromiss er det einaste operative prinsippet.

\subprop{2.2} \textbf{Stilperiodar} er lokale optimum i tilpassingslandskapet: mellombels stabile formkonfigurasjonar som overlever så lenge landskapet ikkje endrar seg. Dei er ikkje årsaker til form, men emergente namn på klynger i formrommet.

\subprop{2.3} Når eit landskap har fleire motstridande mål, oppstår ein \textbf{Pareto-front}. Å formgje er ikkje å finne eit motsetnadslaust optimum, men å velje ein posisjon langs kompromissfronten.

\subprop{2.4} Landskapet er i rørsle. Nye materialar (mahogni, stål, plast) endrar naboskapsstrukturen. Nye teknologiar (dampbøying, sveising, sprøytestøyping) opnar nye regionar. Geopolitiske hendingar (kolonialisme, industrialisering, globalisering) endrar tilpassingsfunksjonen. Kvart slikt skifte er eit \emph{landskapsskifte} som ugyldiggjer eksisterande optimum og tvingar ny navigasjon.


% ================================================================
% PROPOSISJON 3
% ================================================================

\prop{3} Eit logisk bilete av formfakta er at materien ber sin eigen intelligens.

\subprop{3.1} Materialet er ein aktiv, formgjevande og distribuert aktør. Det utgjer eit \textbf{agentialt substrat}. Treverk motstår bøying ut over fiberlengda; stål toler strekk; plast flyt inn i former. Kvar materialeigenskap er ein ibuande seleksjonsregel som avgrensar og mogeleggjer form.

\subprop{3.2} Shannon-entropien for materialar stig frå $H' = 1{,}0$ bits (1200-talet) til $H' = 5{,}07$ bits (1900-talet). Denne kurva er eit kvantitativt mål på det agentiale substratets aukande kompleksitet: jo fleire materialar som er tilgjengelege, jo fleire regionar i formrommet som er navigerbare.

\subprop{3.22} Geometri er materialets ibuande logikk projisert i tre dimensjonar. Forma er ikkje pålagd materialet utanfrå; ho er \emph{forhandla} mellom formgjevar og substrat.

\subprop{3.3} \textbf{Materiell dobbeltheit}, konstruksjonen der berande kjerne er lokalt tre medan fasaden er importert edelt materiale, er eit formuttrykk for det geopolitiske maktsystemet innskriven i sjølve konstruksjonen.


% ================================================================
% PROPOSISJON 4
% ================================================================

\prop{4} Formgjeving er navigasjon utført av kognitive aktørar på eit skalafritt kontinuum.

\subprop{4.1} Navigasjonen skjer på fleire skalaer samstundes: materialet ``bereknar'' kva geometriar det toler; handverkaren ``bereknar'' kva teknikkar som er moglege; marknaden ``bereknar'' kva former som overlever. Ingen enkeltaktør determinerer forma; ho er eit resultat av distribuert forhandling.

\subprop{4.2} Designhistoria vekslar mellom to navigasjonsmodusar:
\begin{enumerate}
  \item \textbf{Utforsking} (\emph{exploration}): raske fasar med stigande entropi, divergerande form, nye materialregime. Døme: 1600-talets materialeksplosjon (+33 materialar), den tidlege modernismen.
  \item \textbf{Utnytting} (\emph{exploitation}): lange fasar med konvergens, fallande entropi, dominans av eit fåtal materialar eller teknikkar. Døme: mahogni-monopolet 1750--1800.
\end{enumerate}

\subprop{4.3} Denne oscillasjonen er isomorf med \emph{punctuated equilibrium} i evolusjonær biologi: lange periodar med stasis avbrote av raske brot. Formhistoria er ikkje lineært framsteg, men ein serie av utforskingar og konsolideringar i eit skiftande landskap.


% ================================================================
% PROPOSISJON 5
% ================================================================

\prop{5} All endeleg formgjeving er frambrakt av eit distribuert svermsinn.

\subprop{5.1} Formgjevinga er aldri ein isolert, deduktiv akt. Ho er eit resultat av samspelet mellom materialets affordanse, verktøyets avgrensing, handverkarens erfaring, og det kulturelle feltet sine seleksjonskriterium. Kvart av desse elementa ber ein form for kognisjon.

\subprop{5.2} Polstring, den mest frekvente teknikken i databasen (618 førekomstar), illustrerer dette: han krev eit berande trevirke \emph{og} eit dekkmateriale \emph{og} ein fyllmasse \emph{og} ein handverkar som meistrar samanføringa. Forma er ikkje vedteken av ein designar; ho er \emph{framforhandla} av eit nettverk av aktørar.

\subprop{5.31} Statistisk maskinlæring viser at stilmerket er ein svak prediktor for eit objekts geometri, medan dimensjonar og materialar predikerer stilperiode med $F1 = 0{,}823$. Dei djupe kreftene (geometri, materiale, teknikk) kodar stil; stilen kodar ikkje geometri.


% ================================================================
% PROPOSISJON 6
% ================================================================

\prop{6} Å formgje er å orkestrere.

\subprop{6.1} Å formgje er å stille opp eit problemrom for eit agentialt materiale, utvide systemets kognitive lyskjegle, opprette dei rette seleksjonsvilkåra, og late det distribuerte nettverket berekne.

\subprop{6.2} Formgjevaren er ikkje ein skapar som deduserer det eine rette svaret. Formgjevaren er ein \emph{navigator} som vel ein posisjon langs ein Pareto-front, ein \emph{kurator} som vel kva seleksjonstrykk som skal dominere, og ein \emph{orkestrator} som set dei agentiale substrata i produktiv spenning.

\subprop{6.3} Funksjonskravet avgrensar formrommet, men determinerer ikkje posisjonen innanfor det. Det som determinerer posisjonen, er tilpassingslandskapet i sin heilskap: materiale, teknologi, økonomi, geografi, kultur.


% ================================================================
% PROPOSISJON 7
% ================================================================

\prop{7} Kvar realisert form er eit mellombels lokalt optimum. Tilpassingslandskapet er i rørsle, og formrommet er prinsipielt ope.

\subprop{7.1} Ingen form er evig optimal. Gårastolen frå 1280 var eit lokalt optimum i eit landskap definert av bjørk, furu og handverkstradisjon. Ein plastkrakk frå 2020 er eit lokalt optimum i eit landskap definert av polypropylen, sprøytestøyping og globalmarknaden. Begge er tilpassa; ingen er universelt ``rette''.

\subprop{7.2} Formrommet utvidar seg irreversibelt. Kvar ny teknikk opnar nye regionar; kvar nytt materiale legg til nye dimensjonar. Dei 65 materialkategoriane og 29 teknikktypane på 1900-talet definerer eit formrom som er ordenar rikare enn dei 2 materiala og 2 teknikkane på 1200-talet.

\subprop{7.21} Dersom nye empiriske data viser at funksjon åleine kan forklare all observert formvariasjon, fell traktaten. Denne falsifiserbarheita er grunnvilkåret for at rammeverket er vitskapeleg.

\subprop{7.3} Å formgje er å setje eit distribuert, agentialt system i vedvarande rørsle mot ein måltilstand som vert tilnærma med aukande presisjon, men aldri endeleg nådd.

\vspace{3em}

\begin{center}
\rule{0.3\textwidth}{0.4pt}\\[1em]
\emph{Fitnesskartet erstattar aksiomet.}
\end{center}

\end{document}
